\setupcolors[state=start]
\setupcolor[hex]
\definecolor[darkblue][h=003366]
\definecolor[darkgray][h=e6e6e6]
\definecolor[redline][h=d9534f]
\definecolor[blueline][h=5bc0de]
\definecolor[greenline][h=5cb85c]

\setuphead[section][style=\bfa, color=darkblue, before=\blank, after=\blank]
\setuphead[subsection][style=\bf, color=black, before=\blank, after=\blank]
\setuppapersize[A4][A4]
\setupbodyfont[11pt]
\setuplayout[width=middle, height=middle, topspace=2cm, bottomspace=2cm]

\starttext
\startsection[title={Quinta Versión: Metaheurística Bio-inspirada (OPP ACO)}]

Además del modelo exacto (OPP\_v3) y la heurística de búsqueda local (Hill Climbing), se diseñó e implementó un motor basado en inteligencia de enjambres: {\bf Ant Colony Optimization (ACO)}. Esta metaheurística iterativa emula el comportamiento de búsqueda de las hormigas reales para converger hacia rutas óptimas.

\startsubsection[title={Definición del Problema Adaptado para ACO}]

La metaheurística ACO, al igual que Hill Climbing, delega la lógica de asignación ({\em order batching}) asumiendo que los lotes ya se encuentran ensamblados. El esfuerzo computacional se concentra exclusivamente en el enrutamiento:

\startitemize[n]
    \item {\bf Input:} Configuración fija de órdenes empaquetadas por lote (parámetro $y_{nb}$) y la topología espacial del almacén ($V, E, d_{ij}$).
    \item {\bf Dinámica Espacial:} Transiciones probabilísticas iterativas de nodos basadas en una matriz de rastros de feromonas (memoria global) y visibilidad heurística (matriz de distancias inversas).
    \item {\bf Objetivo:} Construcción estocástica de secuencias de {\em picking} que minimicen el {\em makespan} global y las penalizaciones temporales $ea_n, ta_n$.
    \item {\bf Output:} Grafo final de trayectorias óptimas compatibles con OPP\_v3.
\stopitemize

\stopsubsection

\startsubsection[title={Fundamentación de la Metaheurística}]

El motor ACO ofrece una propuesta diferencial frente a algoritmos constructivos o de búsqueda local rígida:

\startitemize[bullet]
    \item {\bf Exploración Estocástica:} A diferencia de HC (que suele atraparse en óptimos locales), ACO fomenta la exploración de nuevas ramas topológicas utilizando una matriz de probabilidades de estado de transición.
    \item {\bf Aprendizaje Colectivo:} El rastro de feromonas ($\tau_{ij}$) actúa como un sistema de memoria a largo plazo. Soluciones eficientes encontradas por el enjambre en iteraciones previas sesgan probabilísticamente el camino de futuras hormigas.
    \item {\bf Escalabilidad en Densidad:} Aunque sufre en grafos excesivamente extensos (como se demostró en los test de estrés), ACO es excepcionalmente hábil desenredando marañas de múltiples órdenes dentro de zonas de {\em picking} congestionadas.
\stopitemize

\stopsubsection

\startsubsection[title={Arquitectura Algorítmica y Probabilidad de Transición}]

El núcleo de ACO reside en la ecuación probabilística de selección de estado. Una hormiga en el nodo $i$ elige transitar al nodo $j$ (dentro del conjunto de vecindad no visitada $N_i$) según la siguiente fórmula:

\startformula
P_{ij} = \frac{[\tau_{ij}]^\alpha \cdot [\eta_{ij}]^\beta}{\sum_{k \in N_i} [\tau_{ik}]^\alpha \cdot [\eta_{ik}]^\beta}
\stopformula

Donde:
\startitemize[bullet]
    \item $\tau_{ij}$: Nivel actual de feromona depositada en el arco $(i,j)$. Representa el aprendizaje global.
    \item $\eta_{ij} = \frac{1}{d_{ij}}$: Información heurística local (visibilidad). Favorece los desplazamientos físicamente más cortos.
    \item $\alpha, \beta$: Hiperparámetros que ponderan la importancia relativa del rastro histórico frente a la codicia geográfica.
\stopitemize

\stopsubsection

\startsubsection[title={Actualización de Feromonas (Pheromone Update)}]

Tras finalizar una iteración global (cuando todo el enjambre construyó sus rutas), la matriz de feromonas experimenta dos procesos concurrentes: evaporación y depósito.

\startformula
\tau_{ij} \leftarrow (1 - \rho) \cdot \tau_{ij} + \sum_{k=1}^{m} \Delta \tau_{ij}^k
\stopformula

Donde $\rho$ es el coeficiente de evaporación (evita convergencia prematura), $m$ es el número de hormigas, y $\Delta \tau_{ij}^k$ es la cantidad de feromona depositada por la hormiga $k$, la cual es inversamente proporcional al costo total (función objetivo) de la ruta hallada.

\stopsubsection

\startsubsection[title={Pseudocódigo Estructural}]

\starttyping
procedure ACO_Routing(Bn_fijo, parametros)
  // Inicialización global
  MatrizFeromonas ← Inicializar(\tau_0)
  MejorSolucionGlobal ← Nulo
  MejorCostoGlobal ← \infty
  
  para iteracion = 1 hasta max_iteraciones hacer
    SolucionesHormigas ← []
    
    // Fase constructiva del enjambre
    para k = 1 hasta num_hormigas hacer
      ruta_k ← ConstruirRutaProbabilistica(MatrizFeromonas, Visibilidad)
      costo_k ← EvaluarRuta(ruta_k)
      Agregar ruta_k a SolucionesHormigas
      
      si costo_k < MejorCostoGlobal entonces
        MejorSolucionGlobal ← ruta_k
        MejorCostoGlobal ← costo_k
      fin si
    fin para
    
    // Evaporación global
    Evaporar(MatrizFeromonas, \rho)
    
    // Refuerzo basado en elitismo
    DepositarFeromonas(MatrizFeromonas, SolucionesHormigas)
    
  fin para
  
  devolver MejorSolucionGlobal
fin procedure
\stoptyping

\stopsubsection

\startsubsection[title={Comparativa e Inserción en la Arquitectura OPP}]

La arquitectura bio-inspirada se instanció preservando estrictamente las políticas de DTOs y el patrón Observer descritos en las guías técnicas del proyecto:

\startitemize[bullet]
    \item {\bf OPPAcoModelV3:} Actúa como controlador y hereda abstractamente del orquestador base, compartiendo la misma firma de método `solve(params)`.
    \item {\bf Parámetros Sintonizables:} Además de los hiperparámetros convencionales del OPP ($w_1, w_2, \gamma, \delta$), el modelo exige configuración para el tamaño del enjambre (`ants`), las épocas de búsqueda (`max_iterations`), la evaporación (`evaporation_rate`) y las ponderaciones $\alpha$ y $\beta$.
    \item {\bf Desempeño y Tolerancia:} La estabilización asintótica observada revela que ACO domina la búsqueda local (HC) en instancias altamente laberínticas, pero incurre en un castigo cuadrático computacional al someterse a expansiones territoriales desmedidas (incrementos masivos de $|V|$).
\stopitemize

\stopsubsection

\stopsection
\stoptext